% ---------------------------------------------------------------------
%  CHƯƠNG 1 — MỞ ĐẦU
% ---------------------------------------------------------------------
\chapter{Mở đầu}
\label{ch:modau}

\begin{muctieu}
Chương mở đầu trình bày lý do chọn đề tài, mục tiêu cần đạt, phạm vi thực
hiện và bố cục báo cáo, giúp người đọc nắm được bức tranh tổng quan về đồ án.
\end{muctieu}

\section{Lý do chọn đề tài}

\begin{chuy}
\goiy{nêu bối cảnh thực tế của bài toán, vì sao nó đáng giải quyết, và vì
sao em chọn cách tiếp cận bằng trí tuệ nhân tạo. Có thể dẫn một đến hai số
liệu hoặc một tình huống minh hoạ.}
\end{chuy}

\begin{mau}
Trong bối cảnh thương mại điện tử và mạng xã hội phát triển mạnh, mỗi ngày
có hàng triệu bình luận của người dùng được tạo ra. Việc đọc và tổng hợp
thủ công các bình luận này tốn nhiều công sức và không kịp thời. Bài toán
phân tích cảm xúc \en{(sentiment analysis)} cho phép tự động xác định một
bình luận mang sắc thái tích cực, tiêu cực hay trung tính, từ đó hỗ trợ
doanh nghiệp lắng nghe khách hàng nhanh hơn. Đề tài này được chọn vì vừa có
tính thực tiễn cao, vừa giúp em vận dụng các kiến thức về xử lý ngôn ngữ tự
nhiên và học máy đã tích luỹ.
\end{mau}

\section{Mục tiêu của đồ án}

\begin{chuy}
\goiy{liệt kê ba đến năm mục tiêu cụ thể, đo lường được. Tránh viết chung
chung kiểu "tìm hiểu về AI". Hãy nêu rõ sản phẩm và mức độ hoàn thành mong
muốn, vì Chương~\ref{ch:ketqua} sẽ đối chiếu ngược lại từng mục tiêu này.}
\end{chuy}

\begin{mau}
Đồ án hướng tới các mục tiêu sau:
\begin{itemize}[leftmargin=1.6em, itemsep=2pt]
  \item Xây dựng được một mô hình phân loại cảm xúc cho bình luận tiếng Việt
        với độ chính xác từ 80\% trở lên trên tập kiểm thử.
  \item Cài đặt quy trình tiền xử lý dữ liệu văn bản tiếng Việt hoàn chỉnh.
  \item Phát triển giao diện web cho phép người dùng nhập câu và nhận kết
        quả dự đoán.
  \item So sánh và đánh giá ít nhất hai mô hình học máy khác nhau.
\end{itemize}
\end{mau}

\section{Phạm vi và đối tượng}

\begin{chuy}
\goiy{giới hạn rõ những gì đồ án làm và không làm: loại dữ liệu, ngôn ngữ,
số lượng nhãn, nền tảng triển khai. Giới hạn phạm vi sớm là cách chắc chắn
nhất để kịp hoàn thành trong mười hai tuần.}
\end{chuy}

\begin{mau}
Đồ án tập trung vào phân loại cảm xúc ở mức câu cho văn bản tiếng Việt, với
ba nhãn tích cực, tiêu cực và trung tính. Dữ liệu là bình luận sản phẩm thu
thập công khai. Đồ án không xử lý dữ liệu đa ngôn ngữ, không phát hiện châm
biếm và chưa triển khai ở quy mô sản xuất.
\end{mau}

\begin{luuy}
Lỗi phổ biến nhất của đồ án là chọn phạm vi quá rộng rồi bỏ dở. Một bài toán
hẹp làm tới nơi tới chốn luôn được chấm cao hơn một ý tưởng lớn còn dang dở.
\end{luuy}

\section{Bố cục báo cáo}

\begin{chuy}
\goiy{giữ nguyên đoạn dưới, chỉ sửa nếu em thêm hoặc bớt chương.}
\end{chuy}

Phần còn lại của báo cáo được tổ chức như sau. Chương~\ref{ch:coso} trình
bày cơ sở lý thuyết và công nghệ sử dụng. Chương~\ref{ch:thietke} phân tích
bài toán và thiết kế giải pháp. Chương~\ref{ch:caidat} mô tả quá trình cài
đặt và thử nghiệm. Chương~\ref{ch:ketqua} trình bày kết quả và đánh giá.
Chương~\ref{ch:ketluan} tổng kết và nêu hướng phát triển.
